\documentclass{standalone}
\usepackage{circuitikz}
\usepackage[active,pdftex,tightpage]{preview}
\usepackage{siunitx}
\usepackage[latin1]{inputenc}
\PreviewEnvironment{circuitikz}
\begin{document}

\begin{circuitikz}
    % Define the nodes as coordinates for clarity
    \coordinate (A) at (0, 0);
    \coordinate (B) at (2.5, 0);
    \coordinate (C) at (5, 0);
    \coordinate (D) at (2.5, -3);
    
    

    % Draw the components
    \draw (A) to[R, l=$3$] ++(0,-1.5) to [vsourceAM, l=6] (A |- D)--(D);
    \draw (A) to[R, l=$4$] (B);
    \draw (B) to[R, l=$4$] (C);
    \draw (B) to[R, l=$1$] (D);
    \draw (C) to [vsourceAM, l=$10$] (C |- D)--(D);
    \draw (A)-- ++(0,1)coordinate(Ap) to[R, l=2] (Ap -| C) -- (C);

    %Draw Black dot
    \draw (A) node[circ,label=left:A]{};
    \draw (B) node[circ,label=above:B]{};
    \draw (C) node[circ,label=right:C]{};
    \draw (D) node[circ,label=below:D]{};

    %Units
    \draw (A)++(1,-3.5) node[]{[\si{\volt , \ampere, \ohm}]};

\end{circuitikz}
\end{document}
